\documentclass[11pt]{article}
\usepackage[margin=1in]{geometry}
\usepackage{amsmath,amssymb,amsthm,mathtools}
\usepackage{hyperref}
\usepackage{booktabs}

\newtheorem{theorem}{Theorem}
\newtheorem{lemma}[theorem]{Lemma}
\newtheorem{proposition}[theorem]{Proposition}
\theoremstyle{definition}
\newtheorem{definition}[theorem]{Definition}
\theoremstyle{remark}
\newtheorem*{remark}{Remark}

\newcommand{\Z}{\mathbb{Z}}
\newcommand{\Zq}{\mathbb{Z}_q}
\newcommand{\Rq}{R_q}


\newcommand{\E}{\mathop{\mathbb{E}}}
\newcommand{\lowbits}{\operatorname{LowBits}}
\newcommand{\leak}{\operatorname{leak}}
\newcommand{\geom}{\operatorname{geom}}
\newcommand{\dev}{\operatorname{dev}}
\newcommand{\lean}[1]{\texttt{#1}}
\renewcommand{\Pr}{\mathop{\mathrm{Pr}}}

\title{The order of rejection checks in circl's ML-DSA\\
\large A walk through the Lean proof \lean{RejOrder.lean}}
\author{}\date{}

\begin{document}
\maketitle

\begin{center}
\fbox{\parbox{0.9\textwidth}{\small\textbf{Disclaimer.} This document, the Lean proof and the
numerical experiments were produced mostly by Claude Opus~5.5 (Anthropic), under human
direction. The Lean statements are machine-checked; everything else, including the prose and the
experiments, has not been independently verified.}}
\end{center}

\begin{abstract}
In the signing loop of circl's ML-DSA, the check $\|w_0 - c s_2\|_\infty < \gamma_2-\beta$
(the ``$r_0$ check'') is done \emph{before} the check $\|z\|_\infty < \gamma_1 - \beta$, and both
checks exit early at the first offending coefficient. Timing can therefore reveal, for every
rejected attempt, which check failed and at which coefficient. This note explains a machine-checked
proof (Lean~4 + Mathlib, no \lean{sorry}, standard axioms only) about an attacker who sees
\emph{more} than timing can realistically give: the outcome of every attempt separately, with the
failing check and the exact failing coefficient. Even this view is statistically close to a
distribution that does not depend on the secret: at most $2^{-405}$, $2^{-547}$, $2^{-729}$ in
total variation per attempt for ML-DSA-44/65/87, on average over the public matrix $A$. As a
consequence, for every secret key, all but a $2^{-242}$ fraction of matrices
(ML-DSA-44; $2^{-384}$ for ML-DSA-65, $2^{-566}$ for ML-DSA-87) give an attacker who observes
$2^{35}$ signatures this way an advantage of at most $2^{-128}$. These fractions are crude upper bounds.
We state precisely what is proved, list the modelling assumptions, and sketch the ideas of the
proof.
\end{abstract}

\tableofcontents

%======================================================================
\section{The headline}\label{sec:headline}
%======================================================================

\begin{quote}
\emph{Fix any ML-DSA secret key. For all public matrices $A$ except a small fraction, consider
an attacker who, for each of $2^{35}$ signatures (a thousand per second for a year), observes a
signing attempt in full detail: whether circl's $r_0$ check or $z$ check rejected it, and at
exactly which coefficient. This is far more than a timing measurement reveals. The attacker's
advantage in distinguishing these observations from secret-independent noise is at most
$2^{-128}$.}
\end{quote}

\begin{center}\small
\begin{tabular}{@{}lcc@{}}
\toprule
 & leak per attempt (Thm.~\ref{thm:main}) & fraction of bad $A$ \\
\midrule
ML-DSA-44 & $2^{-405}$ & $\le 2^{-242}$ \\
ML-DSA-65 & $2^{-547}$ & $\le 2^{-384}$ \\
ML-DSA-87 & $2^{-729}$ & $\le 2^{-566}$ \\
\bottomrule
\end{tabular}
\end{center}

\begin{itemize}
  \item The per-attempt bound and the Markov step over $A$ are machine-checked
        (\lean{leak\_mldsa*} and \lean{bad\_frac\_mldsa*}); the arithmetic is in
        Section~\ref{sec:many}.
  \item \textbf{Stronger than timing.} A timing attacker sees one number per signature, its total
        signing time. Our attacker instead sees an attempt's outcome down to the failing
        coefficient. We model one timing measurement as revealing at most one attempt in this much
        detail.
  \item \textbf{One attempt's worth per signature.} Accordingly, the fractions are
        $2^{35}\cdot2^{128}\cdot\varepsilon$: each signature is charged the leak of a single
        attempt. This is a modelling choice, not something we prove. Theorem~\ref{thm:main}
        bounds one attempt. An attacker who saw \emph{every} attempt of a signature in full
        detail would be charged for every attempt, $\E[N]\cdot\varepsilon$ with $N$ the number of
        attempts ($\E[N]\approx4.3$, $5.1$, $3.9$). Charging $\E[N]$ instead would make each
        fraction $2$--$3$ bits larger ($2^{-239}$, $2^{-381}$, $2^{-564}$).
  \item \textbf{These fractions are crude upper bounds.} Bad matrices are probably much rarer.
        Markov's inequality uses nothing but the average over $A$, which is the worst possible
        case for it: all the leak concentrated on a few matrices. The average is itself
        overestimated: it is dominated by pairs of nonces with many degenerate NTT slots, which
        the proof charges the trivial bound (Section~\ref{sec:proof}). And the bound charges for
        information a timing attacker never sees: which coefficient failed. Numerical experiments put
        the true per-attempt deviation for a random $A$ thousands of bits below the proved bound
        (Section~\ref{sec:tight}).
  \item The trade-off is linear: a weaker advantage target $2^{-128+d}$ makes the fraction $2^d$
        times smaller. For example, for ML-DSA-44, advantage at most $2^{-64}$ holds for all but a
        $2^{-306}$ fraction of $A$.
  \item Not covered: the later $ct_0$ and hint checks, and timing variation from other parts of
        signing. The hybrid
        argument over attempts and the random-oracle modelling of the challenge are standard but
        not formalised. See Section~\ref{sec:assumptions}.
\end{itemize}

%======================================================================
\section{Setting}
%======================================================================

\paragraph{What circl does.}
One iteration of the rejection loop (\texttt{sign/mldsa/*/internal/dilithium.go}) is:
\begin{enumerate}
  \item sample $y \in \Rq^L$ with coefficients in $(-\gamma_1,\gamma_1]$;
  \item $w = Ay$, $(w_1,w_0) = \mathrm{Decompose}(w)$, $\tilde c = H(\mu \,\|\, w_1)$, $c = \mathrm{SampleInBall}(\tilde c)$;
  \item \textbf{$r_0$ check:} if $\|w_0 - cs_2\|_\infty \ge \gamma_2-\beta$, reject;
  \item \textbf{$z$ check:} $z = y + cs_1$; if $\|z\|_\infty \ge \gamma_1-\beta$, reject;
  \item further checks ($\|ct_0\|_\infty \ge \gamma_2$, hint weight $>\omega$), then output.
\end{enumerate}
FIPS~204 performs the $z$ check first. Both \texttt{Exceeds} routines exit at the first failing
polynomial, and (in the generic code) at the first failing coefficient; the AVX2 code exits at the
first failing block of 32 coefficients. So the timing of steps 3--4 reveals (at most) the index of
the first failing coefficient of the first failing check.

\paragraph{Why the order could matter.}
The $z$ check alone is harmless: for every shift $u = cs_1$ with $\|u\|_\infty\le\beta$, each
coordinate of $y+u$ lands outside $(-(\gamma_1-\beta), \gamma_1-\beta)$ with probability exactly
$p_Z = (2\beta+1)/2\gamma_1$, independently, so the first failing index is \emph{exactly} geometric
and independent of the secret. The $r_0$ check is different: its outcome depends on $v = cs_2$
through $\lowbits(Ay)-v$, and the coordinates of $Ay$ are not independent of each other. If it is
done first, its failure index is revealed for every $y$, not only for those that already passed the
$z$ check. The question is whether that index carries information about $v$.

%======================================================================
\section{The model}
%======================================================================

Throughout, $m = Kn$ and $l = Ln$ with $n=256$, and vectors in $\Rq^L$, $\Rq^K$ are identified
with $\Z^{l}$ or $\Zq^{l}$, $\Zq^m$ coefficient-wise, polynomial by polynomial (index
$256\,i + s$ for coefficient $s$ of polynomial $i$). The challenge $c$ is \emph{fixed}; it enters only
through the two shifts
\[
  u = cs_1 \in \Z^{l}, \qquad v = cs_2 \in \Z^{m}, \qquad \|u\|_\infty,\ \|v\|_\infty \le \beta,
\]
which are otherwise arbitrary (\lean{Bounded}). We need $\beta = \tau\eta$, which
bounds $\|cs_i\|_\infty$ for every $c$ of weight $\tau$.

\begin{definition}[Parameters, \lean{Params}]
Naturals $q,\gamma_1,\gamma_2,\beta,b$ with $q = 2\gamma_2 b + 1$, $\beta<\gamma_1$,
$\beta<\gamma_2$, $b>0$. (Here $b = (q-1)/2\gamma_2$ is the number of high-bit values.)
\end{definition}

\begin{center}
\begin{tabular}{lcccccc}
\toprule
 & $(K,L)$ & $\gamma_1$ & $\gamma_2$ & $\beta$ & $b$ \\
\midrule
ML-DSA-44 & $(4,4)$ & $2^{17}$ & $95232$  & $78$  & $44$ \\
ML-DSA-65 & $(6,5)$ & $2^{19}$ & $261888$ & $196$ & $16$ \\
ML-DSA-87 & $(8,7)$ & $2^{19}$ & $261888$ & $120$ & $16$ \\
\bottomrule
\end{tabular}
\end{center}
with $q = 8380417$ throughout.

\begin{definition}[Model of one attempt]\label{def:model}
\leavevmode
\begin{itemize}
  \item $\mathcal{B} = (-\gamma_1,\gamma_1]^{l} \subset \Z^l$; $y$ is uniform on $\mathcal{B}$ (\lean{box}).
  \item $\lowbits:\Zq\to\Z$ is the low part of FIPS~204 Decompose with $\alpha = 2\gamma_2$,
        including the corner case $r - r_0 = q-1 \Rightarrow r_0 \gets r_0 - 1$ (\lean{lowBits}).
  \item $M \in \Zq^{m\times l}$ is the matrix of $A$; $w = M\,(y \bmod q)$.
  \item Coordinate $i$ fails the $r_0$ check iff $|\lowbits(w_i) - v_i| \ge \gamma_2-\beta$;
        coordinate $k$ fails the $z$ check iff $|y_k + u_k| \ge \gamma_1-\beta$.
  \item $I(y)\in\{0,\dots,m-1\}\cup\{\bot\}$ and $J(y)\in\{0,\dots,l-1\}\cup\{\bot\}$ are the
        least failing indices.
  \item The \emph{view} is
  \[
    \mathrm{view}_{M,u,v}(y) =
    \begin{cases}
      \mathsf{rFail}(i) & \text{if } I(y) = i,\\
      \mathsf{zFail}(j) & \text{if } I(y) = \bot,\ J(y)=j,\\
      \mathsf{pass}(y+u) & \text{if } I(y)=J(y)=\bot.
    \end{cases}
  \]
\end{itemize}
\end{definition}

\paragraph{Faithfulness to circl.}
The circl code computes $w_0 - cs_2$ in $\Zq$ and compares its centred representative with
$\gamma_2-\beta$. Since $|w_0|\le\gamma_2$ and $|v_i|\le\beta$, the integer $w_0-v_i$ is well
inside $(-q/2,q/2)$, so this is exactly $|\lowbits(w_i)-v_i|$. The same holds for $z$. The
coefficient ordering, the ranges of $y$, the parameters and the matrix $A$ as a block matrix of
negacyclic multiplications all match. A coarser view (e.g.\ AVX2's block index) is a function of
the modelled view and is therefore covered.

\begin{definition}[Reference distribution and leak, \lean{leak}]
Let $p_R = 1 - b\,(2(\gamma_2-\beta)-1)/q$ and $p_Z = (2\beta+1)/2\gamma_1$; these are the
per-coordinate failure probabilities of the $r_0$ check for \emph{uniform} $w\in\Zq$ and of the
$z$ check. Let $\geom_p(i) = (1-p)^i p$ and $\geom_p(\bot)=(1-p)^{N}$ for $N$ trials. Define
\[
  \leak(M,u,v) = \sum_{i<m} \Bigl|\Pr_y[\mathsf{rFail}(i)] - \geom_{p_R}(i)\Bigr|
  + \sum_{j<l} \Bigl|\Pr_y[\mathsf{zFail}(j)] - \geom_{p_Z}(j)\,(1-p_R)^m\Bigr|.
\]
The reference values do not depend on $u$, $v$ or $M$.
\end{definition}

%======================================================================
\section{The statements}
%======================================================================

\begin{theorem}[\lean{leak\_mldsa44}, \lean{leak\_mldsa65}, \lean{leak\_mldsa87}]\label{thm:main}
Let $(q,\gamma_1,\gamma_2,\beta,b,K,L)$ be an ML-DSA parameter set, $\Rq = \Zq[X]/(X^{256}+1)$,
and $A$ uniform in $\Rq^{K\times L}$. For every $u\in\Z^{Ln}$, $v\in\Z^{Kn}$ with
$\|u\|_\infty,\|v\|_\infty\le\beta$,
\[
  \E_A\bigl[\leak(A,u,v)\bigr] \le
  \begin{cases}
    2^{-405} & \text{ML-DSA-44},\\
    2^{-547} & \text{ML-DSA-65},\\
    2^{-729} & \text{ML-DSA-87}.
  \end{cases}
\]
\end{theorem}

The accepted views are handled exactly:

\begin{proposition}[\lean{view\_pass}]\label{prop:pass}
For every $M$, bounded $u,v$ and $z\in\Z^l$,
\[
  \Pr_y[\mathrm{view}(y) = \mathsf{pass}(z)] =
  \begin{cases}
    (2\gamma_1)^{-l} & \text{if } \|z\|_\infty < \gamma_1-\beta \text{ and}\\
       & \quad \forall i,\ |\lowbits\bigl((Mz - (Mu+v))_i\bigr)| < \gamma_2-\beta,\\
    0 & \text{otherwise.}
  \end{cases}
\]
In particular it depends on $(u,v)$ only through $Mu+v$, which for ML-DSA is $c\,t$ with
$t = As_1+s_2$.
\end{proposition}

The proof uses the shift lemma (Lemma~3 of the Dilithium specification, \lean{lowBits\_shift}):
for $|v|\le\beta$, $|\lowbits(w)-v| < \gamma_2-\beta \iff |\lowbits(w-v)|<\gamma_2-\beta$.

Together these give the meaning of $\leak$:

\begin{theorem}[\lean{indist}]\label{thm:indist}
Fix $M$. Let $(u,v)$ and $(u',v')$ be bounded with $Mu+v = Mu'+v'$. Then for every test
$D$ on views,
\[
  \Bigl|\Pr_y[D(\mathrm{view}_{M,u,v}(y))] - \Pr_y[D(\mathrm{view}_{M,u',v'}(y))]\Bigr|
  \le \leak(M,u,v)+\leak(M,u',v').
\]
\end{theorem}

\begin{proof}[Proof idea]
Write both probabilities as sums over views. The $\mathsf{pass}$ terms are equal by
Proposition~\ref{prop:pass}. Each $\mathsf{rFail}$/$\mathsf{zFail}$ term differs by at most its
distance to the common reference value, summed over the two sides.
\end{proof}

\paragraph{How to read these.}
Theorem~\ref{thm:main} says that the rejection part of the view of one attempt is, on average over
$A$, within $\varepsilon$ (in total variation) of a fixed distribution that does not depend on the
secret. The accepted part reveals $z$, which is in the signature anyway, with the exact
distribution of Proposition~\ref{prop:pass}. Theorem~\ref{thm:indist} is a clean corollary but
assumes the same $c\,t$ (the full $t$, not only $t_1$); the useful content is that $\leak$ is small.

\subsection{Many signatures}\label{sec:many}

An attacker watches many attempts. For fixed shifts, the views of independent attempts are
within the \emph{sum} of the per-attempt leaks of a secret-independent reference (a standard
hybrid argument). With random challenges the shifts are random, which we cover by weighting the
attempts. Markov's inequality over $A$ then gives:

\begin{theorem}[\lean{bad\_frac}, \lean{bad\_frac\_mldsa44/65/87}]\label{thm:headline}
Let $\varepsilon$ be the bound of Theorem~\ref{thm:main}. For every finite family of bounded
shifts $(u_i,v_i)$ with weights $w_i\ge0$ of total weight $W=\sum_i w_i$, and every $\eta>0$,
\[
  \Pr_A\Bigl[\textstyle\sum_i w_i\,\leak(A,u_i,v_i)\ \ge\ \eta\Bigr] \;\le\; \frac{W\varepsilon}{\eta}.
\]
\end{theorem}

For the headline of Section~\ref{sec:headline}: $2^{35}$ signatures, each charged as one
attempt, with that attempt's weight spread uniformly over the challenges it could receive, so
$W=2^{35}$; take $\eta=2^{-128}$. Then the fraction of bad $A$ is at most
$2^{35+128}\varepsilon$, i.e.\ $2^{-242}$, $2^{-384}$, $2^{-566}$. Charging every attempt instead
means giving attempt $i$ of a signature weight $\Pr[N\ge i]$, so $W=2^{35}\,\E[N]$.

What is \emph{not} in Lean is the hybrid step from ``sum of leaks'' to ``attacker advantage over
many attempts'', and the random-oracle step that makes the challenges uniform and independent of
$y$ (assumption A1 below). Both are standard. The bad set of matrices depends on the key.

%======================================================================
\section{Assumptions and gaps}\label{sec:assumptions}
%======================================================================

\begin{enumerate}\renewcommand{\labelenumi}{\textbf{A\arabic{enumi}.}}
\item \textbf{Fixed challenge.} In circl, $c$ is derived from $H(\mu\,\|\,w_1)$ with
  $w_1 = \mathrm{HighBits}(Ay)$, so $c$ depends on $y$. The model fixes $c$ and takes $y$ uniform.
  This is justified in the random-oracle model: for a fresh query, $c$ is uniform and independent
  of $y$, and since Theorem~\ref{thm:main} holds for \emph{every} bounded $(u,v)$, it also holds
  averaged over $c$. Consequently, leaking $c$ (e.g.\ through \textsc{SampleInBall} timing) is
  harmless for this argument; leaking anything about $w_1$ or $y$ beyond the view is not covered.
\item \textbf{Uniform $y$ and $A$.} $y$ (from SHAKE-256) and $A$ (from \textsc{ExpandA}) are
  modelled as uniform, and $A$ as independent of $(s_1,s_2)$.
\item \textbf{Average over $A$.} The bound is on $\E_A[\leak]$, not for every $A$. Markov's
  inequality turns it into a per-key statement; see Theorem~\ref{thm:headline}.
\item \textbf{One attempt.} Theorem~\ref{thm:main} covers one attempt. Theorem~\ref{thm:headline}
  bounds weighted sums of leaks over many attempts, but the hybrid argument that turns this sum
  into a bound on the attacker's advantage is not formalised. The headline charges one attempt
  per signature (Section~\ref{sec:headline}); a proof would charge $\E[N]$, costing $2$--$3$
  bits.
\item \textbf{Only the $r_0$ and $z$ checks.} The view stops at $\mathsf{pass}(z)$. The
  later $\|ct_0\|_\infty\ge\gamma_2$ check (also early-exiting) and the hint-weight check are not
  modelled. They are independent of the order question but are timing-visible.
\item \textbf{Leakage model.} Each \texttt{Exceeds} call is assumed to reveal only the first
  failing index; all other operations (NTT, hashing, \textsc{Decompose}, \ldots) are assumed
  constant time.
\end{enumerate}

%======================================================================
\section{Proof overview}\label{sec:proof}
%======================================================================

The proof has five steps.
\begin{enumerate}
  \item \emph{Exact counting:} the $z$ check is exactly geometric; the $r_0$ check is exactly
        geometric when $w$ is replaced by a uniform vector.
  \item \emph{Product sets:} every rejection view is a difference of events ``$Ay$ lies in a
        product set'', so the leak is bounded by deviations of such events.
  \item \emph{Fourier:} averaged over $A$, the square of each deviation is a sum over pairs of nonces
        $(y,y')$ of covariances, which Fourier analysis bounds by a weighted count of test vectors
        orthogonal to the image of $A \mapsto (Ay, Ay')$.
  \item \emph{Module structure:} for $\Rq = \Zq[X]/(X^n+1)$ those test vectors live on the few
        NTT slots where $y$ and $y'$ are linearly dependent, and they are counted row by row.
  \item \emph{Numerics:} one integer inequality per parameter set, checked by the kernel.
\end{enumerate}

%----------------------------------------------------------------------
\subsection{Step 1: exact counting}
%----------------------------------------------------------------------

\begin{lemma}[\lean{z\_pass\_count}, \lean{J\_dist}]
For $|u_k|\le\beta$, exactly $2(\gamma_1-\beta)-1$ points $y_k\in(-\gamma_1,\gamma_1]$ pass the $z$
check. Since $\mathcal{B}$ is a product, $\Pr_y[J(y) = o] = \geom_{p_Z}(o)$ for every bounded $u$.
\end{lemma}

\begin{lemma}[\lean{r\_pass\_count}]
For $|v_i|\le\beta$, exactly $b\,(2(\gamma_2-\beta)-1)$ elements $w\in\Zq$ satisfy
$|\lowbits(w)-v_i|<\gamma_2-\beta$. Hence for $w$ uniform in $\Zq^m$, the first failing index is
distributed as $\geom_{p_R}$, and $\Pr_w[\text{all } i\in T \text{ fail}] = p_R^{|T|}$.
\end{lemma}

The only nontrivial ingredient is a bijection between $\Zq$ and the valid decompositions
$(r_1,r_0)$: $\{0,\dots,b-1\}\times(-\gamma_2,\gamma_2]$ together with the corner point
$(0,-\gamma_2)$, which has exactly $q$ elements. On this set $\lowbits$ is the projection to $r_0$,
so shifting by $v_i$ shifts $r_0$, and the count reduces to counting an interval.

%----------------------------------------------------------------------
\subsection{Step 2: views as product sets}
%----------------------------------------------------------------------

Let $\mathrm{good}_r \subseteq \Zq$ be the values of coordinate $r$ of $w$ that pass the $r_0$
check (with shift $v_r$), and let
\[
  P_i = \prod_{r<i}\mathrm{good}_r \times \prod_{r\ge i}\Zq
\]
be the set of $w$ whose first $i$ coordinates pass (\lean{Pset}), of density
$\mu_i = (1-p_R)^i$ (\lean{muP\_eq}). Then ``the $r_0$ check first fails at $i$'' is
$Ay \in P_i\setminus P_{i+1}$, and ``the $z$ check first fails at $j$'' is $y \in B_j$ and
$Ay \in P_m$, where $B_j = \{y\in\mathcal B : J(y) = j\}$ has probability $\geom_{p_Z}(j)$
exactly (Step~1). For $B'\subseteq\mathcal{B}$ let
\[
  D_i(B') = \Bigl|\Pr_y\bigl[y\in B',\ Ay \in P_i\bigr] - \Pr_y[y\in B']\,\mu_i\Bigr|
  \qquad(\text{\lean{Dev}}).
\]
Comparing each view with its reference value term by term gives (\lean{leak\_le\_Dev})
\[
  \leak(A,u,v) \le \sum_{i<m}\bigl(D_i(\mathcal B) + D_{i+1}(\mathcal B)\bigr)
    + \sum_{j<l} D_m(B_j).
\]
So it suffices to bound $\E_A D_i(B')$, uniformly in $i$ and in the sub-event $B'$.

%----------------------------------------------------------------------
\subsection{Step 3: the second moment by Fourier analysis}
%----------------------------------------------------------------------

Expanding the square,
\[
  \E_A D_i(B')^2 = \frac{1}{|\mathcal B|^2}\sum_{y,y'\in B'}
    \mathrm{Cov}_A\bigl(\mathbf 1[Ay\in P_i],\,\mathbf 1[Ay'\in P_i]\bigr),
\]
a sum over pairs of nonces of covariances over $A$. For a fixed pair, $A\mapsto(Ay,Ay')$ is an
additive map, and $P_i$ is a \emph{product} set. Its indicator has Fourier coefficients that factor
over coordinates. A nonzero frequency on a free coordinate ($r \ge i$) gives $0$, and on a
passing coordinate costs a factor of at most $b = p_R/(1-p_R) \approx 2^{-10.25}$ relative to the
zero frequency, since the bad set is small. Averaging the characters over $A$ keeps exactly the
pairs of frequencies $(\xi,\xi')$ for which $\xi\cdot Ay + \xi'\cdot Ay' = 0$ for all $A$. This gives
(\lean{cov\_le})
\[
  \bigl|\mathrm{Cov}_A\bigr| \le \mu_i^2\!\!\sum_{\substack{\xi,\xi'\ne0\\
    \xi\cdot Ay+\xi'\cdot Ay'\equiv0}}\!\! b^{\mathrm{wt}(\xi)+\mathrm{wt}(\xi')},
\]
with $\mathrm{wt}$ the Hamming weight. If no nonzero test vector $(\xi,\xi')$ is orthogonal to
the image, $(Ay,Ay')$ is uniform and the covariance is $0$; otherwise the bound charges each
orthogonal test vector by its weight. The rest of the proof counts these test vectors.

%----------------------------------------------------------------------
\subsection{Step 4: the module structure}
%----------------------------------------------------------------------

Up to here nothing used the ring. Now $A$ is a $K\times L$ matrix of elements of
$\Rq = \Zq[X]/(X^n+1)$, and the only property used (\lean{IsNTTMul}) is:

\begin{quote}
multiplication is diagonalised by evaluation at $n$ distinct nonzero points $\zeta_k$, i.e.\
$\widehat{ac}(k) = \hat a(k)\,\hat c(k)$ where $\hat a(k) = \sum_s a_s\zeta_k^s$.
\end{quote}

\noindent For ML-DSA, $\zeta_k = 1753^{2k+1}$ (proved: $1753^{256} = -1$ and $1753$ has order
$512$; \lean{isNTTMul\_mldsa}).

\noindent In Lean, Theorem~\ref{thm:main} is stated over $\Rq$ itself: section \lean{RingForm}
takes $\Rq$ = \lean{AdjoinRoot} $(X^{256}+1)$ from Mathlib, $A$ uniform in
$\Rq^{K\times L}$, and $M = $ \lean{ringMat}$\,A$, the $\Zq$-matrix of $y \mapsto Ay$ on
coefficient vectors (\lean{ringMat\_mulVec}). The bridge to the negacyclic matrices is exact: a
product in $\Rq$ has the coefficients of the negacyclic product (\lean{coeffs\_mul}, by
evaluating at the $\zeta_k$ and inverting the Vandermonde matrix), and a uniform $A$ has uniform
coefficients (\lean{avg\_ringMat}).

\paragraph{Dual coordinates and degenerate slots.}
Write a test vector $\xi$ on row $i$ of the output as
$\xi_{(i,s)} = \sum_k \zeta_k^s\,\tilde\xi_i(k)$, which is always possible since the
Vandermonde matrix is invertible (\lean{dual}). Then
$\xi\cdot(Ay) = \sum_{i,j,k}\tilde\xi_i(k)\,\hat a_{ij}(k)\,\hat y_j(k)$, and since the
$\hat a_{ij}(k)$ can be chosen freely, $(\xi,\xi')$ is orthogonal for all $A$ exactly when
\[
  \tilde\xi_i(k)\,\hat y_j(k) + \tilde\xi'_i(k)\,\hat y'_j(k) = 0
  \qquad\text{for all } i,j,k
\]
(\lean{slot\_of\_orth}). Call slot $k$ \emph{degenerate} for the pair if the vectors
$(\hat y_j(k))_j$ and $(\hat y'_j(k))_j$ in $\Zq^L$ are linearly dependent, and let $D(y,y')$ be
the set of degenerate slots (\lean{degS}). An orthogonal pair of test vectors has all its dual
coordinates on $D(y,y')$ (\lean{dualIn\_of\_orth}), and if $D = \emptyset$ there is none
(\lean{dualIn\_empty}).

\paragraph{Degenerate slots are rare (window argument).}
\begin{lemma}[\lean{slot\_prob}, \lean{deg\_card\_prob}]
If $2\gamma_1\le q$, then for $y,y'$ independent and uniform on $\mathcal{B}$, a given set of $d$
slots is degenerate with probability at most $\rho^d$, $\rho = (q+1)/(2\gamma_1)^L$.
\end{lemma}
\begin{proof}[Idea]
A union bound over one of the $q+1$ lines $(c:c')$ per slot reduces this to one linear condition
on $(\hat y_j(k),\hat y'_j(k))$ per slot and per $j<L$. Say $r$ slots constrain $y_j$. Fix all
coefficients of $y_j$ except the first $r$. The remaining unknowns form a polynomial of degree $<r$
whose values at $r$ distinct points are prescribed, so it is determined (\lean{window}). Because
$2\gamma_1\le q$, coefficients in $(-\gamma_1,\gamma_1]$ are determined by their residues mod $q$.
So at most a $(2\gamma_1)^{-r}$ fraction of the box survives, for each of the $L$ polynomials.
\end{proof}

\paragraph{Counting test vectors by rows.}
Let $E_b(d)$ bound the $b$-weighted number of nonzero test vectors whose duals lie on a given set
$S$ of $d$ slots, $\sum_{\xi\ne0,\ \tilde\xi\subseteq S} b^{\mathrm{wt}(\xi)}$
(\lean{wsum\_le}). Group the $\xi$ by their coefficient support $T$. On row $i$, $\xi$ vanishes
outside the $t_i$ positions of $T$ in that row, so on a cyclic run of $g_i$ consecutive positions
(\lean{cgap}; runs may wrap around because $\zeta_k^n=-1$ for all $k$, \lean{dual\_row\_zero}).
A function $\tilde\xi_i$ supported on at most $r$ slots whose ``power sums'' vanish on $r$
consecutive exponents is zero (\lean{dual\_vanish}, an uncertainty principle), so $\tilde\xi_i$ is
determined by its values on all but $g_i$ slots of $S$, and row $i$ contributes at most
$q^{\min(t_i,\,d-g_i)}$ test vectors (\lean{card\_suppIn1\_le\_xg}). The weight
$b^{|T|}$ factors over the rows too, so the sum over $T$ is the $K$-th power of a row sum
(\lean{wcount\_le\_rows}). The row sum groups the subsets of an $n$-cycle by their longest free
run; those whose free runs all have length at most $g$ are counted through compositions, as
coefficients of $(1+X+\dots+X^g)^{v}$ (\lean{card\_linSet}, \lean{card\_shortSet}), evaluated by
packing polynomial coefficients into large integers (Kronecker substitution, \lean{GapEval}).

\paragraph{Putting it together.}
Group the pairs $(y,y')$ by $d = |D(y,y')|$. Pairs with $d = 0$ contribute nothing. For $1\le
d\le d_m$, the pair sum is at most $E_b(d)^2$ ($\xi$ and $\xi'$ both live on $D$), and the
probability of $d$ given degenerate slots is at most $\rho^d$. Pairs with more than $d_m$
degenerate slots are charged the trivial bound. So (\lean{Dsq\_le})
\[
  \E_A D_i(B')^2 \le \sum_{d=1}^{d_m}\binom nd\rho^d E_b(d)^2 + \binom{n}{d_m+1}\rho^{d_m+1},
\]
uniformly in $i$ and $B'$ (the pairs run over $B'\times B'\subseteq\mathcal B\times\mathcal B$).
Jensen then gives $\E_A D_i(B')$ as the square root (\lean{avg\_Dev\_le}).

%----------------------------------------------------------------------
\subsection{Step 5: numerics}
%----------------------------------------------------------------------

Cleared of denominators, ``the right-hand side above is at most $2^{-2j}$'' is one inequality
between natural numbers per parameter set (\lean{numR}), checked by kernel evaluation
(\lean{decide +kernel}, not \lean{native\_decide}; \lean{numR\_44} etc.). Then
$\E_A[\leak] \le (2Kn + Ln)\,2^{-j}$ (\lean{leak\_avg\_fourier}).

\begin{center}
\begin{tabular}{lcccc}
\toprule
 & $d_m$ & $j$ & $p_R$ & bound on $\E_A[\leak]$ \\
\midrule
ML-DSA-44 & 19 & 417 & $8.2\cdot10^{-4}$ & $2^{-405}$ \\
ML-DSA-65 & 20 & 560 & $7.5\cdot10^{-4}$ & $2^{-547}$ \\
ML-DSA-87 & 23 & 742 & $4.6\cdot10^{-4}$ & $2^{-729}$ \\
\bottomrule
\end{tabular}
\end{center}
The bound is limited by the pairs with more than $d_m$ degenerate slots and by the crude count of
test vectors on many slots, not by actual leakage (Section~\ref{sec:tight}).

\paragraph{History.} An earlier version of this proof used Bonferroni truncation over the events
``all coordinates in $T$ fail'', a capped second moment and counts by the size of $T$, and gave
$2^{-289}$, $2^{-418}$, $2^{-568}$. The product-Fourier bound, developed for the nonce-reuse
analysis in \lean{Bathwater.lean}, replaced it; the old route is in the git history.

%======================================================================
\section{How tight is this? (numerical experiments)}\label{sec:tight}
%======================================================================

The proved bounds are limited by the method, not by actual leakage. This section collects
evidence for that. Nothing here is formalised; the experiments are numerical, and the
extrapolations are labelled as such.

\paragraph{Two easy facts.}
\begin{itemize}
  \item The leak is never exactly zero. $\Pr_y[I(y)=0]$ is a fraction with denominator
        $(2\gamma_1)^l$, a power of two, while $p_R$ has the odd prime $q$ in its denominator. So
        $\leak(A,u,v)\ge 1/(q(2\gamma_1)^l)\approx2^{-18455}$ for ML-DSA-44. This is distance to
        the reference, not dependence on the secret.
  \item Some matrices really leak, so averaging over $A$ is unavoidable. If the first row of $A$
        is a constant polynomial $c$ and zero elsewhere, then $w_0 = c\,y_0$, and for a suitable
        $c$ whether coordinate $0$ fails depends visibly on $v_0$, at a level of about
        $1/(2\gamma_1)$. Such matrices make up only about $q^{-1023}$ of all $A$.
\end{itemize}

\paragraph{One coordinate, computed exactly.}
For one output coordinate $w_0=\sum_c M_{0c}y_c$ and a uniformly random row of $M$, the shift
lemma makes ``coordinate $0$ fails, given $v_0$'' a translate of a fixed set. Its dependence on
$v_0$ is a Fourier sum over the $q-1$ nonzero frequencies $\xi$, with
\[
  \hat P(\xi)=\prod_{c}\varphi(M_{0c}\,\xi),\qquad
  |\varphi(x)|=\frac{|\sin(2\pi\gamma_1x/q)|}{2\gamma_1|\sin(\pi x/q)|}.
\]
We evaluated this sum exactly, in $\log_2$, for all $\xi$. The method was checked against brute
force on small instances.

\begin{center}
\begin{tabular}{lccc}
\toprule
 & deviation of one coordinate & proved per-attempt bound & gap \\
\midrule
ML-DSA-44 & $\approx2^{-18100}$ & $2^{-405}$ & ${\approx}17700$ bits \\
ML-DSA-65 & $\approx2^{-25250}$ & $2^{-547}$ & ${\approx}24700$ bits \\
ML-DSA-87 & $\approx2^{-35400}$ & $2^{-729}$ & ${\approx}34700$ bits \\
\bottomrule
\end{tabular}
\end{center}
(Several random rows per parameter set; these are typical values, varying by tens of bits.) The
scale is easy to explain. Since $x\mapsto M_{0c}x$ permutes $\Zq$, each factor $\log_2|\varphi|$
averages exactly $-\log_2(2\gamma_1)$, so a typical frequency gives $2^{-18\cdot1024}$ for
ML-DSA-44. The best of the roughly $2^{22}$ frequencies gains about $300$ bits through ordinary
fluctuations, not because a few factors are unusually large.

\paragraph{Pairs of coordinates, sampled.}
A pair has about $q^2\approx2^{46}$ frequencies, too many to scan; that would take about five
years on 24 cores. Instead, for ML-DSA-44 we sampled $4\cdot10^9$ random frequency pairs and
scanned about $6000$ lines $\xi_2=\lambda\xi_1$ exactly. Findings:
\begin{itemize}
  \item At a fixed frequency the terms are exactly independent when $A$ is uniform. The sampled
        distribution of $\log_2|\hat P|$ matches the resulting model to sampling precision, in
        moments and in the tail out to about $6\sigma$.
  \item The same extrapolation applied to single lines, whose exact maxima we know, predicts
        them correctly: on average to within $0.2$ bits, with the right spread.
  \item The extrapolated maximum for a pair is about $2^{-17979\pm10}$, about $150$ bits larger
        than for one coordinate.
  \item No structured direction stands out: not the axes, small $\lambda$, simple fractions, or,
        for two adjacent coefficients of one polynomial, the directions singled out by the
        negacyclic structure. One sampled pair showed a mild excess in the extreme tail
        ($p\approx0.04$), worth at most $5$--$10$ bits if real.
\end{itemize}

\paragraph{Larger joint events.}
Under the same model, joint events with all of $T$ failing, for $|T|$ up to about $30$--$50$,
are still estimated near $2^{-16000}$ for
ML-DSA-44. Very large joint events are genuinely non-uniform: for ML-DSA-65 and 87, $m>l$, so
some nonzero frequency gives exactly $1$. But these enter with weight about $p_R^{|T|}$, roughly
$2^{-10|T|}$, so they do not change the picture.

\paragraph{Conclusion.}
The true per-attempt leak for a random $A$ is plausibly in the range $2^{-16000}$ to
$2^{-18000}$ for ML-DSA-44, far below the proved $2^{-405}$, and the fractions of bad matrices
in the headline are correspondingly pessimistic. Closing the gap with a proof would need a quite
different technique.

%======================================================================
\section{Map of the Lean file}
%======================================================================

\begin{center}
\begin{tabular}{ll}
\toprule
Section & Content \\
\midrule
\lean{Defs} & the model (Definition~\ref{def:model}) \\
\lean{Layer0} & Decompose bijection, shift lemma, exact counts, $J$ geometric \\
\lean{Layer1} & accepted views (Proposition~\ref{prop:pass}) \\
\lean{Leak} & $\leak$, \lean{indist}, averages, Markov \\
\lean{Uncertainty} & Vandermonde window and uncertainty principle \\
\lean{Module}, \lean{ModuleProb} & block matrix, dual coordinates, slot conditions, window argument \\
\lean{ModuleCount}, \lean{Profile} & test vectors on a coordinate set; \lean{deg\_card\_prob} \\
\lean{Negacyclic} & $\Zq[X]/(X^n+1)$ satisfies \lean{IsNTTMul}; ML-DSA roots \\
\lean{RingForm} & $\Rq$ as Mathlib's \lean{AdjoinRoot}; \lean{coeffs\_mul}, \lean{ringMat}, \lean{avg\_ringMat} \\
\lean{Numeric} & binomial coefficients for the kernel \\
\lean{GapCount}, \lean{GapComb} & cyclic free runs; counting sets by free runs \\
\lean{GapEval}, \lean{GapSum} & Kronecker substitution; regrouping the sum over rows \\
\lean{GapNumeric} & helpers for the kernel-checked sums \\
\lean{Covariance} & Fourier bound on covariances of product-set indicators (\lean{cov\_le}) \\
\lean{TestVectors} & weighted counts of test vectors (\lean{wsum\_le}, \lean{wcount\_le\_rows}) \\
\lean{FourierLeak} & \lean{leak\_le\_Dev}, \lean{Dsq\_le}, \lean{avg\_Dev\_le}, \lean{leak\_avg\_fourier} \\
\lean{MLDSA} & Theorem~\ref{thm:main} (kernel checks \lean{numR\_44} etc.) \\
\lean{Headline} & many attempts, few bad matrices (Theorem~\ref{thm:headline}) \\
\bottomrule
\end{tabular}
\end{center}

The file builds with Lean~4.34 and a recent Mathlib in about 20 seconds.
\lean{\#print axioms} reports only \lean{propext}, \lean{Classical.choice} and \lean{Quot.sound}
for \lean{indist}, the three main theorems and the three \lean{bad\_frac} instances.

\end{document}
